\section{Introduction}
\label{sec:intro}

Vision transformers have rapidly become a cornerstone in modern computer vision, achieving state-of-the-art results in tasks ranging from image classification to object detection and segmentation \cite{he2022masked} \cite{oquab2023dinov2} \cite{radford2021clip} \cite{touvron2022deitiii}. These models leverage the transformer architecture to process images as sequences of patches. With their ability to model long-range dependencies and capture global context, vision transformers \cite{dosovitskiy2021vit} have demonstrated remarkable performance across a wide variety of benchmarks \cite{chen2015coco} \cite{deng2009imagenet} \cite{pascal-voc-2012} \cite{young2014flickr} \cite{zhou2017ade}.

The unique designs of vision transformers have given rise to intriguing behaviors that are not yet fully understood. One such phenomenon is the emergence of a subset of tokens that exhibit exceptionally high activation norms in certain layers of the network. These tokens, which we refer to as ``\textit{massive tokens}'' (occasionally abbreviated as ``MA''), appear to dominate the attention distribution, effectively acting as ``\textit{attention sinks}'' \cite{gu2024attnsink} that influence the overall flow of information through the network.

In addition to massive tokens, our investigations reveal the presence of what we term ``\textit{artifact tokens}.'' These tokens do not naturally exhibit the extreme activation norms of massive tokens; however, they become evident under specific conditions---taking on the extreme-magnitude and attention-sink characterization of massive tokens only when the original massive tokens have been masked or removed. This observation suggests that vision transformers possess a built-in redundancy mechanism, where a limited number of tokens are capable of assuming the role of massive tokens if needed. 

These observations carry both theoretical interest and practical implications. We demonstrate that by strategically leveraging the distinct roles of massive and artifact tokens, it is possible to reconfigure the model’s attention dynamics to improve computational efficiency and boost performance.

In light of these observations, our work primarily makes the following contributions:
\begin{itemize}
    \item Fast Nystr\"om Attention (FNA), a training-free method that approximates self-attention at inference in linear time and space complexity using the properties of massive and artifact tokens.
    \vspace{-0.2cm}
    \item A novel, training-free algorithm that efficiently identifies massive tokens in vision transformers, enabling extraction and strategic masking with minimal computational overhead during inference. We show that our proposed masking procedure yields consistent performance improvements across a range of downstream multi-modal and dense prediction tasks.
    \vspace{-0.2cm}
    \item A comprehensive analysis of the mechanisms in vision transformers that facilitate the formation of massive tokens and their correlation with artifact tokens.
\end{itemize}

